% Part: history
% Chapter: biographies
% Section: rozsa-peter

\documentclass[../../../include/open-logic-section]{subfiles}

\begin{document}

\olfileid[tr]{his}{bio}{pet} 

\olsection{R\'ozsa P\'eter}

\olphoto{peter-rozsa}{R\'ozsa P\'eter}
 
R\'ozsa P\'eter, R\'ozsa Politzer adıyla 17 Şubat 1905'te Macaristan'ın
Budapeşte kentinde doğdu. En çok, özyineleme teorisi alanının kurulmasında
temel önemde olan özyinelemeli fonksiyonlar üzerine çalışmalarıyla tanınır.

P\'eter ağır siyasal koşullarda büyüdü---ergenlik yıllarında I. Dünya Savaşı
sürüyordu---ama Budapeşte'deki seçkin Maria Terezia Kız Okulu'na gidebildi ve
1922'de buradan mezun oldu. Ardından Budapeşte'deki P\'azm\'any P\'eter
Üniversitesi'nde (daha sonra Lor\'and E\"otv\"os Üniversitesi adını aldı)
öğrenim gördü. Babasının ısrarıyla kimya okumaya başladı; daha sonra
matematiğe geçti ve 1927'de mezun oldu. Lise matematiği öğretmek için gerekli
yeterliliğe sahip olsa da Büyük Buhran dünya ekonomisini etkilediğinden
dönemin ekonomik durumu vahimdi. P\'eter bu sırada özel öğretmenlik ve bire
bir matematik dersleri gibi geçici işlerde çalıştı. Sonunda matematikte
lisansüstü öğrenim görmek üzere üniversiteye döndü. Başlangıçta sayı teorisi
üzerine çalışmayı tasarlamıştı; fakat vardığı sonuçların daha önce ispatlanmış
olduğunu öğrenince matematiği neredeyse tümüyle bırakıyordu. G\"odel'in
eksiklik teoremleri üzerinde çalışmaya teşvik edildi ve onun bazı sonuçlarını
farkında olmadan farklı yollarla ispatladı. Bu, özgüvenini geri kazandırdı;
P\'eter, David Hilbert'in temeller programından esinlenen ilk özyineleme
teorisi makalelerini yazmayı sürdürdü. Doktorasını 1935'te aldı ve 1937'de
\emph{Journal of Symbolic Logic} dergisinin editörü oldu.

P\'eter'in ilk makaleleri, özyinelemeli fonksiyon teorisi alanının kurucu
katkıları olarak yaygın biçimde kabul edilir. \cite{Peter1935a}'da farklı
özyineleme türleri arasındaki ilişkiyi inceledi. \cite{Peter1935b}'de belirli
bir özyinelemeli tanımlı fonksiyonun ilkel özyinelemeli olmadığını gösterdi.
Bu, Wilhelm Ackermann'a ait daha önceki bir sonucu basitleştirdi. P\'eter'in
basitleştirdiği fonksiyon bugün çoğunlukla Ackermann fonksiyonu---bazen daha
yerinde bir adlandırmayla Ackermann--P\'eter fonksiyonu---diye bilinir.
Özyinelemeli fonksiyon teorisi üzerine ilk kitabı yazdı \citep{Peter1951}.

Çalışmalarının önemine ve etkisine karşın P\'eter, 1945'e kadar tam zamanlı
bir öğretim görevi alamadı. II. Dünya Savaşı sırasında Macaristan'ın Nazi
işgali altında olduğu dönemde Yahudi karşıtı yasalar nedeniyle P\'eter'in
ders vermesine izin verilmedi. Hükûmet 1944'te Budapeşte'de bir Yahudi gettosu
kurdu; getto
kentin geri kalanından tecrit edilmişti ve silahlı muhafızlarca korunuyordu.
P\'eter, getto kurtarılana kadar, 1945'e dek burada yaşamak zorunda kaldı.
Ardından Budapeşte Öğretmen Yetiştirme Yüksekokulu'nda, 1955'ten başlayarak da
E\"otv\"os Lor\'and Üniversitesi'nde ders verdi. Matematikte Akademi Doktoru
olan ilk Macar kadın matematikçi ve Macar Bilimler Akademisi'ne seçilen ilk
kadın oldu.

P\'{e}ter, yalnızca ders anlatmak yerine öğrencileriyle matematiksel
problemlerin doğasını ve güzelliğini araştırmayı yeğleyen tutkulu bir
matematik öğretmeni olarak tanınırdı. Bu nedenle öğrencileri ona sevgiyle
``Rosa Teyze'' derdi. P\'eter 1977'de 71 yaşında öldü.

\begin{reading}
Daha fazla biyografik okuma için \citep{Oconnor2014} ve
\citep{Andrasfai1986} kaynaklarına bakın. \citet{Tamassy1994}, P\'eter'le kısa
bir söyleşi yapmıştır. Matematik üzerine eğlenceli bir okuma için P\'eter'in
\emph{Playing With Infinity} adlı kitabına bakın \citep{Peter2010}.
\end{reading}

\end{document}
